\newpage
\appendix
\section{Training Details}
\label{app:details}

\paragraph{SFT stage.}
We fine-tune Qwen3-32B using LoRA with rank 64 and alpha 128 on all
linear layers. The SFT dataset consists of 5{,}000 samples: 637 Chinese
mathematical critique examples and 4{,}363 English math reasoning problems
from NuminaMath-CoT and MetaMathQA. Training runs for 3 epochs (120
steps) with learning rate $5 \times 10^{-5}$, cosine decay with 3\%
warmup, batch size 2, gradient accumulation 8, and max sequence length
8{,}192. Final SFT loss: 0.311, token accuracy: 89.7\%.

\paragraph{GRPO stage.}
We train with the TopoPRM composite reward on 637 Chinese mathematical
critique prompts. Each prompt generates 8 completions with temperature
0.7, top-$p$ 0.95, and max completion length 2{,}048. We adopt the
MS-Swift \emph{server mode} architecture, where a dedicated vLLM rollout
server (4 GPUs, tensor-parallel) handles high-throughput generation while
4 separate GPUs perform LoRA gradient updates. After each optimisation
step, updated adapter weights are synchronised back to the server via
NCCL broadcast, ensuring the generation policy stays current. Training
runs for 1 epoch with learning rate $5 \times 10^{-6}$, cosine decay,
per-device batch size 2, and gradient accumulation 4.

\paragraph{Reward weights.}
Default weights: $w_o = 0.4$ (outcome), $w_t = 0.2$ (topology),
$w_c = 0.15$ (continuity), $w_f = 0.15$ (format), $w_l = 0.1$ (length).

\paragraph{Hardware.}
All experiments run on 8$\times$ NVIDIA L20X GPUs (143\,GB VRAM each).
For GRPO, we partition GPUs into an inference group (4 GPUs running the
vLLM rollout server with tensor parallelism) and a training group
(4 GPUs running distributed LoRA optimisation), following the
MS-Swift~\citep{msswift} server-mode paradigm.

\section{SCAE Algorithm}
\label{app:scae}

Algorithm~\ref{alg:scae} presents the full procedure of Stratified
Clipping Advantage Estimation.

\begin{algorithm}[H]
\caption{Stratified Clipping Advantage Estimation (SCAE)}
\label{alg:scae}
\begin{algorithmic}[1]
\Require Group of completions $\{c^{(1)}, \ldots, c^{(G)}\}$ for prompt $q$
\Require Outcome rewards $\{R_\mathrm{out}^{(i)}\}_{i=1}^G$, auxiliary rewards $\{r_\mathrm{aux}^{(i)}\}_{i=1}^G$
\State $\bar{r}_\mathrm{acc} \leftarrow \frac{1}{G}\sum_{i=1}^G R_\mathrm{out}^{(i)}$
\State $\mathcal{C} \leftarrow \{i : R_\mathrm{out}^{(i)} = 1\}$ \Comment{correct set}
\State $\mathcal{W} \leftarrow \{j : R_\mathrm{out}^{(j)} = 0\}$ \Comment{incorrect set}
\State $\bar{r}_\mathrm{aux}^+ \leftarrow \mathrm{mean}(\{r_\mathrm{aux}^{(i)} : i \in \mathcal{C}\})$ if $\mathcal{C} \neq \emptyset$ else $0$
\State $\bar{r}_\mathrm{aux}^- \leftarrow \mathrm{mean}(\{r_\mathrm{aux}^{(j)} : j \in \mathcal{W}\})$ if $\mathcal{W} \neq \emptyset$ else $0$
\For{$i \in \mathcal{C}$}
    \State $A^{(i)} \leftarrow (1 - \bar{r}_\mathrm{acc}) + \max(0,\; r_\mathrm{aux}^{(i)} - \bar{r}_\mathrm{aux}^+)$
\EndFor
\For{$j \in \mathcal{W}$}
    \State $A^{(j)} \leftarrow -\bar{r}_\mathrm{acc} + \min(0,\; r_\mathrm{aux}^{(j)} - \bar{r}_\mathrm{aux}^-)$
\EndFor
\Return Advantages $\{A^{(i)}\}_{i=1}^G$
\end{algorithmic}
\end{algorithm}

\section{DAG Extraction Details}
\label{app:dag}

\paragraph{Step segmentation.}
We identify reasoning step boundaries using the following markers:
\begin{itemize}[leftmargin=*,itemsep=1pt]
    \item Numbered steps: patterns such as ``Step 1:'', ``(1)'', and Chinese-style step markers   \item Mathematical delimiters: ``$\because$'', ``$\therefore$'',
          ``$\Rightarrow$''
    \item Logical connectives: ``hence'', ``therefore'', and their
          Chinese equivalents
    \item Line breaks with significant indentation changes
\end{itemize}

\paragraph{Dependency detection.}
For each pair of steps $(r_i, r_j)$ with $i < j$, we check:
\begin{enumerate}[leftmargin=*,itemsep=1pt]
    \item \textbf{Expression overlap}: whether $r_j$ contains
    mathematical expressions (extracted via regex) that first appeared
    in $r_i$.
    \item \textbf{Variable reference}: whether $r_j$ uses variables or
    intermediate results defined in $r_i$.
    \item \textbf{Sequential proximity}: consecutive steps
    $(r_i, r_{i+1})$ receive weak sequential edges (weight 0.3).
\end{enumerate}

\paragraph{Step classification.}
Each node is classified by keyword matching into one of seven types:
\textsc{definition}, \textsc{derivation}, \textsc{computation},
\textsc{conclusion}, \textsc{auxiliary}, \textsc{substitution}, and
\textsc{case\_analysis}.

\section{Reward Function Details}
\label{app:rewards}

\paragraph{Outcome reward.}
Parses the model's \texttt{<answer>} JSON to extract the predicted
student score and critique conclusion. Exact score match yields 1.0;
within $\pm 1$ yields 0.5. Matching the conclusion adds 0.5 bonus.
The total is capped and normalised to $[0, 1]$.

\paragraph{Format reward.}
Checks for \texttt{<think>$\cdots$</think>} and
\texttt{<answer>$\cdots$</answer>} blocks with valid JSON:
both present $\to$ 1.0; only valid \texttt{<answer>} $\to$ 0.3;
otherwise $\to$ 0.0.

\paragraph{Length reward.}
Linear decay from 1.0 (at $\leq$ 2{,}000 characters) to 0.0 (at
$\geq$ 4{,}000 characters).

\paragraph{Topological structure reward.}
Defined in \S\ref{sec:method}.
The topology reward decomposes into six auditable components with default weights:
valid DAG ($\lambda_b{=}0.20$), acyclicity ($\lambda_a{=}0.15$),
no-orphan conclusions ($\lambda_o{=}0.15$), directional consistency
($\lambda_d{=}0.15$), step alignment ($0.10$), and reference-edge F1
($\lambda_k{=}0.25$, redistributed when no reference graph is available).
A hard gate zeros the topology score when the extracted graph is invalid.

\paragraph{Dynamic weight adaptation.}
The composite reward supports optional dynamic weight adaptation.
At each batch, per-component correlation with outcome correctness and
spread (standard deviation) are computed.
A reliability score $\rho_m = 0.5 \cdot \mathrm{corr}(r_m, r_o) + 0.5 \cdot \min(1, 2\sigma_m)$
modulates each weight, blended with the static baseline via
$\eta{=}0.50$.
An outcome-floor constraint ($w_o \geq 0.35$) is enforced after
renormalisation to prevent correctness signal erosion.

\paragraph{Continuity reward.}
Defined in \S\ref{sec:method}. Uses regex-based expression matching
to verify that each step references expressions or claims from prior
steps or the problem statement.
Full coverage ($\eta{=}1.0$) yields reward 1.0; partial coverage yields
$0.8\,\eta$.

\section{Limitations}
\label{app:limitations}

Three limitations of the current work warrant discussion. First, the DAG
extraction relies on rule-based parsing of inter-step references through
keyword matching and expression overlap. This procedure may miss implicit
logical dependencies that lack explicit surface-level cues, such as
unstated algebraic identities or geometric properties applied without
reference. Second, the continuity reward uses regex-based expression
matching rather than semantic understanding; it detects when a step
references an expression from a previous step but cannot identify subtle
logical gaps or incorrect applications of mathematical principles that
reuse the same symbols. Third, our evaluation focuses exclusively on
mathematical reasoning in a Chinese educational critique setting. Whether
TopoPRM generalises to other structured reasoning domains, such as code
generation or scientific reasoning, remains an open question.

\section{Future Work}
\label{app:future}

Three directions follow from this work. First, replacing the rule-based
DAG extractor with a lightweight learned parser could capture implicit
dependencies while remaining fast enough for online reward computation
during GRPO training. Second, combining the verifiable structural rewards
of TopoPRM with learned PRMs~\citep{lightman2023lets} could provide
complementary supervision: structural rewards enforce global coherence
while learned rewards capture fine-grained step correctness. Third,
extending the framework to multi-modal mathematical reasoning, where
diagrams and figures introduce spatial dependency structures beyond text,
would test the generality of topology-aware rewards in richer settings.

\section{Case Study}
\label{app:case_study}

We present a representative example comparing the reasoning structure
of completions from the SFT baseline, outcome-only GRPO, and TopoPRM
on a high-school algebra problem requiring multi-step equation solving
with verification of intermediate steps.

\paragraph{SFT output.}
The SFT model produces a correct answer but with redundant intermediate
steps and one orphan conclusion that does not connect to the main
derivation chain. The reasoning DAG contains a back-edge from the
conclusion to an earlier auxiliary step, violating acyclicity.

\paragraph{Outcome-only GRPO output.}
The outcome-only model reaches the correct answer more concisely but
skips a key substitution step, producing a reasoning trace with a
continuity gap. The topological structure reward assigns a low continuity
score to this trace.

\paragraph{TopoPRM output.}
The TopoPRM model produces a concise, well-structured derivation where
each step follows logically from its predecessors. The extracted DAG is
acyclic, fully connected, and directionally consistent. The reasoning
trace is 23\% shorter than the SFT output while maintaining correctness.

\section{Distillation Details}
\label{app:distill}

For reasoning compression, we use the TopoPRM-trained Qwen3-32B
(GRPO checkpoint-79) as the teacher model to generate high-quality
reasoning traces. We filter teacher outputs to retain only those with
composite reward $R_\mathrm{total} > 0.7$. The student model (Qwen3-8B)
is trained using reverse-KL distillation with LoRA (rank 32, alpha 64)
for 2 epochs, max sequence length 2{,}048, learning rate
$3 \times 10^{-5}$, and a maximum of 4{,}000 training samples.
On public benchmarks, the distilled student achieves 82.5\% on GSM8K
(1{,}319 problems) and 59.8\% on MATH-500.

\section{Additional Results}
\label{app:additional}
% Training curves, per-category breakdowns, and additional ablations
% will be added after all experiments complete.
